\documentclass[11pt]{article}

\usepackage{xeCJK}
\setCJKmainfont{STSong}
\setCJKsansfont{Heiti SC}

\input{../../templates/template.LaTeX.homework/structure.tex}

%----------------------------------------------------------------------------------------
%  ASSIGNMENT INFORMATION
%----------------------------------------------------------------------------------------
\newcommand{\assignmentQuestionName}{习题}
\newcommand{\assignmentClass}{离散数学}
\newcommand{\assignmentTitle}{第3章 关系 — 第9次作业}
\newcommand{\assignmentAuthorName}{乐绎华}
\newcommand{\studentID}{23363017}
\newcommand{\assignmentClassInstructor}{左孝凌、刘永才《离散数学》}
\newcommand{\assignmentDueDate}{}

\begin{document}

\maketitle
\thispagestyle{empty}
\newpage

%----------------------------------------------------------------------------------------
%  3-9(1) 集合划分数
%----------------------------------------------------------------------------------------

\begin{question}{3-9(1) 集合划分数}
  4个元素的集合共有多少个不同的划分？

  \begin{answer}
    4个元素的集合的划分数是 Bell 数 $B_4 = 15$。

    \vspace{0.5em}
    \textbf{详细推导（按分块数分类）：}

    \begin{itemize}
      \item \textbf{1块}：$\{\{a,b,c,d\}\}$ — 1种
      \item \textbf{2块}：有两种情况
        \begin{itemize}
          \item $\{1,3\}$ 型：选1个元素单独一组，$\binom{4}{1}=4$种
          \item $\{2,2\}$ 型：$\binom{4}{2}/2=3$种（选2个元素为第一块，其余为第二块，但块无序）
        \end{itemize}
        小计：$4+3=7$种
      \item \textbf{3块}：$\{1,1,2\}$ 型：选哪2个元素合在一组，$\binom{4}{2}=6$种
      \item \textbf{4块}：$\{\{a\},\{b\},\{c\},\{d\}\}$ — 1种
    \end{itemize}

    \vspace{0.5em}
    \textbf{总计}：$1 + 7 + 6 + 1 = 15$种
  \end{answer}
\end{question}

%----------------------------------------------------------------------------------------
%  3-10(1) 等价关系的并集
%----------------------------------------------------------------------------------------

\begin{question}{3-10(1) 等价关系的并集}
  设 $R$ 和 $R'$ 是集合 $A$ 上的等价关系，用例子证明 $R \cup R'$ 不一定是等价关系。

  \begin{answer}
    \textbf{反例}：设 $A = \{1, 2, 3\}$

    取 $R = \{\langle 1,1\rangle, \langle 2,2\rangle, \langle 3,3\rangle, \langle 1,2\rangle, \langle 2,1\rangle\}$（自反、对称、传递）

    取 $R' = \{\langle 1,1\rangle, \langle 2,2\rangle, \langle 3,3\rangle, \langle 2,3\rangle, \langle 3,2\rangle\}$（自反、对称、传递）

    则
    $$R \cup R' = \{\langle 1,1\rangle, \langle 2,2\rangle, \langle 3,3\rangle, \langle 1,2\rangle, \langle 2,1\rangle, \langle 2,3\rangle, \langle 3,2\rangle\}$$

    \vspace{0.5em}
    \textbf{问题}：$\langle 1,2\rangle \in R \cup R'$，$\langle 2,3\rangle \in R \cup R'$，但 $\langle 1,3\rangle \notin R \cup R'$，故\textbf{传递性不成立}。

    因此 $R \cup R'$ 不一定是等价关系。
  \end{answer}
\end{question}

%----------------------------------------------------------------------------------------
%  3-10(2) 4元素集合的等价关系数
%----------------------------------------------------------------------------------------

\begin{question}{3-10(2) 4元素集合的等价关系数}
  由4个元素组成的有限集上所有的等价关系的个数为多少？

  \begin{answer}
    等价关系与集合的划分一一对应。4个元素的集合有15种不同的划分（见 3-9(1)），故有 \textbf{15个等价关系}。

    \begin{info}[说明:]
      这是因为：给定一个划分，定义关系 $R$ 为"在同一分块中"，则 $R$ 是等价关系；反之，给定一个等价关系，其等价类构成一个划分。二者一一对应。
    \end{info}
  \end{answer}
\end{question}

%----------------------------------------------------------------------------------------
%  3-10(6) 对称传递关系的自反性
%----------------------------------------------------------------------------------------

\begin{question}{3-10(6) 对称传递关系的自反性}
  设 $R$ 是集合 $A$ 上的对称和传递关系，证明如果对于 $A$ 中的每一个元素 $a$，在 $A$ 中同时也存在一个 $b$，使 $\langle a,b\rangle \in R$，则 $R$ 是一个等价关系。

  \begin{answer}
    已知 $R$ 是对称的和传递的，只需证明\textbf{自反性}。

    \vspace{0.5em}
    设 $a \in A$，由题设存在 $b \in A$ 使 $\langle a,b\rangle \in R$。

    由对称性：$\langle b,a\rangle \in R$

    由传递性：$\langle a,b\rangle \in R$ 且 $\langle b,a\rangle \in R$ 蕴含 $\langle a,a\rangle \in R$

    因此对所有 $a \in A$，$\langle a,a\rangle \in R$，即 $R$ 是自反的。

    \vspace{0.5em}
    故 $R$ 是等价关系（自反、对称、传递）。$\square$
  \end{answer}
\end{question}

%----------------------------------------------------------------------------------------
%  3-10(8) 复数集合上的等价关系
%----------------------------------------------------------------------------------------

\begin{question}{3-10(8) 复数集合上的等价关系}
  设 $C$ 是实数部分非零的全体复数组成的集合，$C^*$ 上关系 $R$ 定义为：$(a+bi)\ R\ (c+di) \Leftrightarrow ac > 0$，证明 $R$ 是等价关系，并给出 $R$ 的等价类的几何说明。

  \begin{answer}
    \vspace{1em}
    \noindent\textbf{1. 自反性}：对任意 $z = a+bi \in C^*$，$a \neq 0$，故 $a \cdot a = a^2 > 0$，所以 $zRz$。

    \vspace{1em}
    \noindent\textbf{2. 对称性}：若 $z_1 = a+bi$，$z_2 = c+di$，$z_1Rz_2$，则 $ac > 0$。由于乘法可交换，$ca > 0$，故 $z_2Rz_1$。

    \vspace{1em}
    \noindent\textbf{3. 传递性}：若 $z_1Rz_2$ 且 $z_2Rz_3$，设 $z_1 = a+bi$，$z_2 = c+di$，$z_3 = e+fi$。

    由 $z_1Rz_2$：$ac > 0$

    由 $z_2Rz_3$：$ce > 0$

    若 $a > 0$ 则 $c > 0$（因为 $ac > 0$），从而 $e > 0$（因为 $ce > 0$），故 $ae > 0$

    若 $a < 0$ 则 $c < 0$，从而 $e < 0$，故 $ae > 0$

    因此 $z_1Rz_3$。

    \vspace{1em}
    \noindent\textbf{等价类的几何说明}：

    $R$ 将复平面分割成两类：
    \begin{itemize}
      \item \textbf{第一类}：实部为正的所有复数（不包括纯虚数，即 $a > 0$ 的复数）
      \item \textbf{第二类}：实部为负的所有复数（不包括纯虚数，即 $a < 0$ 的复数）
    \end{itemize}

    这是因为 $ac > 0$ 意味着 $a$ 和 $c$ 同号。几何上，两个复数有等价关系当且仅当它们位于复平面同一半平面（不包括虚轴）。

    \begin{info}[几何意义:]
      复平面被虚轴分成左右两个半平面。等价关系 $R$ 判定两个实部非零的复数是否落在同一个半平面内。
    \end{info}
  \end{answer}
\end{question}

\end{document}
